\section{Preliminaries}

\label{sec:prelim}

\paragraph{Discrete diffusion models}
Let $\dv=(x_1,\dots,x_L)\in\Vcal^L$ denote a discrete sequence over a finite vocabulary $\Vcal$, and let $p_0(\dv)$ denote the target data distribution.
A diffusion model defines generation through a learned backward process,
\begin{equation}
    p_\theta\rbr{\dv_{0:T}} = \textstyle p_T\rbr{\dv_T}\prod_{t=1}^{T}p_{\theta,t}\rbr{\dv_{t-1}\mid \dv_t},
    \label{eq:backward_joint}
\end{equation}
where $\dv_T$ is sampled from a simple initial distribution, $\dv_0$ denotes target samples, and $\dv_{1:T-1}$ are auxiliary intermediate variables. 
Without additional structure, these intermediate variables are latent, making it generally intractable to directly maximize the marginal likelihood for $p_\theta\rbr{\xb_0}$. Conventional diffusion models bypass this complication by prescribing a forward corruption process,
\begin{equation}
    q\rbr{\dv_{0:T}} = \textstyle p\rbr{\xb_0}\prod_{t=1}^{T}q_t\rbr{\dv_t\mid \dv_{t-1}},
    \label{eq:forward_joint}
\end{equation}
which defines an auxiliary joint distribution over clean data and corrupted states. 
Training can then be written as joint distribution matching, 
\begin{equation}
    \min_{\theta}\,\,\KL\!\rbr{
    q\rbr{\dv_{0:T}}\,\|\, p_\theta\rbr{\dv_{0:T}}
    }.
    \label{eq:joint_kl}
\end{equation}
For diffusion language modeling, a common choice is masked diffusion models~(MDM), where the forward process $q_t$ is absorbing-state masking corruption~\citep{MD4, MDLM},
\begin{equation}
    q_t\rbr{z_i\mid x_i}=\alpha_t \mathbf{1}\{z_i=x_i\}+ (1-\alpha_t)\mathbf{1}\{z_i=\masktoken\}, 
    \label{eq:mdm_forward}
\end{equation}
where each token is either kept unchanged or replaced by a mask token $\masktoken$ controlled by some predefined masking schedule $\alpha_t$.
The resulting masked diffusion objective reduces to a weighted denoising loss over masking-corrupted positions, 
\begin{equation}
    \Lcal_{\mathrm{MDM}}\rbr{\theta} =\textstyle
    \EE_{t,\dv_0,\dv_t\sim q_t(\cdot\mid \dv_0)}
    \sbr{\lambda_t\sum_{i:\dv_{t,i}=\masktoken}-\log p_\theta\rbr{\dv_{0,i}\mid \dv_t,t}
    },
    \label{eq:mdm_objective}
\end{equation}
where $\lambda_t$ is determined by the masking schedule. 

\paragraph{Complications in forward process design}
The objective in Eq.~\eqref{eq:mdm_objective} highlights the central role of the forward corruption process, which specifies both the training inputs $\dv_t$ and the distributional path $q\rbr{\dv_{0:T}}$ that the learned backward process learns to invert. 
Thus, the difficulty of learning and inference is strongly determined by how the forward process decomposes the gap between the initial distribution and the target distribution. 
While Gaussian perturbations provide a natural local geometry in continuous diffusion~\citep{ddpm,score-based}, discrete language modeling has no intrinsic analogue. 
For MDM with absorbing-state masking in Eq.~\eqref{eq:mdm_forward}, the induced neighborhood of token $a$ degenerates to $\Ncal_t(a)=\cbr{a,\masktoken}$, where each token can only remain unchanged or collapse to the shared absorbing state. Therefore, no direct token-to-token proximity is preserved~\citep{SCUD, mdpo}. 
As a result, the model is trained on mathematically convenient intermediate states that may be misaligned with the drafts and errors encountered during generation. 
